\hypertarget{error-handling-and-robustness}{%
\chapter{ERROR HANDLING AND
ROBUSTNESS}\label{error-handling-and-robustness}}

\hypertarget{error-handling-robustness}{%
\chapter{ERROR HANDLING \& ROBUSTNESS}\label{error-handling-robustness}}

\hypertarget{error-handling-debugging-c98}{%
\chapter{ERROR HANDLING \& DEBUGGING
(C++98)}\label{error-handling-debugging-c98}}

\hypertarget{assertions}{%
\section{1. ASSERTIONS}\label{assertions}}

Use \passthrough{\lstinline!assert!} to catch programming logic errors
during development.

\begin{lstlisting}[language={C++}]
#include <iostream>
#include <cassert>

int divide(int a, int b) {
    assert(b != 0);  // Terminates program if b is 0
    return a / b;
}

int main() {
    std::cout << divide(10, 2) << "\n";
    return 0;
}
\end{lstlisting}

\begin{center}\rule{0.5\linewidth}{0.5pt}\end{center}

\hypertarget{exceptions}{%
\section{2. EXCEPTIONS}\label{exceptions}}

C++ uses exceptions to handle runtime errors gracefully.

\hypertarget{basic-try-catch}{%
\subsection{2.1 Basic Try-Catch}\label{basic-try-catch}}

\begin{lstlisting}[language={C++}]
#include <iostream>

int safe_divide(int a, int b) {
    if (b == 0) {
        throw "Division by zero condition!";
    }
    return a / b;
}

int main() {
    try {
        int result = safe_divide(10, 0);
        std::cout << result << "\n";
    } catch (const char* msg) {
        std::cerr << "Error: " << msg << "\n";
    }
    return 0;
}
\end{lstlisting}

\hypertarget{catching-multiple-types}{%
\subsection{2.2 Catching Multiple Types}\label{catching-multiple-types}}

\begin{lstlisting}[language={C++}]
try {
    // code...
} catch (int e) {
    std::cerr << "Integer exception: " << e << "\n";
} catch (const char* e) {
    std::cerr << "String exception: " << e << "\n";
} catch (...) {
    std::cerr << "Unknown exception caught!\n";
}
\end{lstlisting}

\hypertarget{standard-exceptions}{%
\subsection{2.3 Standard Exceptions}\label{standard-exceptions}}

Inherit from \passthrough{\lstinline!std::exception!} for custom
exceptions.

\begin{lstlisting}[language={C++}]
#include <iostream>
#include <exception>
#include <stdexcept> // runtime_error, logic_error

class MyError : public std::exception {
public:
    virtual const char* what() const throw() {
        return "My Custom Error";
    }
};

void test() {
    throw std::runtime_error("Standard Runtime Error");
}

int main() {
    try {
        test();
    } catch (const std::exception& e) {
        std::cerr << "Caught: " << e.what() << "\n";
    }
    return 0;
}
\end{lstlisting}

\begin{longtable}[]{@{}
  >{\raggedright\arraybackslash}p{(\columnwidth - 0\tabcolsep) * \real{0.06}}@{}}
\toprule
\endhead
\#\#\# Professional Notes: Debugging \& Quality Assurance \\
\#\#\#\# 1. Unit Testing in C++ Testing is essential for maintaining
large codebases. Modern C++ typically uses third-party frameworks: *
\textbf{Google Test (GTest)}: The industry standard. Uses macros like
\passthrough{\lstinline!EXPECT\_EQ!},
\passthrough{\lstinline!ASSERT\_TRUE!}. * \textbf{Catch2}: Header-only,
very easy to integrate. Uses natural BDD style:
\passthrough{\lstinline!REQUIRE(x == 42)!}. \\
\#\#\#\# 2. Debugging with GDB and LLDB Command-line debuggers allow you
to inspect the program state at runtime. *
\textbf{\passthrough{\lstinline!break [file]:[line]!}}: Set a
breakpoint. * \textbf{\passthrough{\lstinline!print [var]!}}: Inspect
variable value. * \textbf{\passthrough{\lstinline!backtrace!} (bt)}:
Show the current call stack. *
\textbf{\passthrough{\lstinline!watch [var]!}}: Pause whenever a
variable's value changes. \\
\#\#\#\# 3. Defensive Programming Techniques * \textbf{Static Assertions
(C++11)}:
\passthrough{\lstinline!static\_assert(sizeof(int) == 4, "32-bit int required");!}.
Checks conditions at compile time. *
\textbf{\passthrough{\lstinline!noexcept!}}: Mark functions that are
guaranteed not to throw. Allows the compiler to generate more optimized
code and skip stack unwinding logic. * \textbf{Core Dumps}: On Linux,
enable core dumps (\passthrough{\lstinline!ulimit -c unlimited!}) to
capture the memory state of a crashed program for post-mortem
analysis. \\
\bottomrule
\end{longtable}

\hypertarget{exception-specifications-c98}{%
\subsection{2.4 Exception Specifications
(C++98)}\label{exception-specifications-c98}}

In C++98, you can specify what a function might throw. (Note: Deprecated
in C++11, but valid here).

\begin{lstlisting}[language={C++}]
// Can throw only int
void func() throw(int) {
    throw 42;
}

// Cannot throw anything
void no_throw() throw() {
    // logic...
}
\end{lstlisting}

\hypertarget{stack-unwinding-raii}{%
\subsection{2.5 Stack Unwinding \& RAII}\label{stack-unwinding-raii}}

When an exception is thrown, the stack is unwound, identifying
destructors for all local objects.

\begin{lstlisting}[language={C++}]
class Resource {
public:
    Resource() { std::cout << "Acquired\n"; }
    ~Resource() { std::cout << "Released\n"; }
};

void risky() {
    Resource r; // Destructor guaranteed to run
    throw 1;
}

int main() {
    try {
        risky();
    } catch (...) {
        std::cout << "Caught\n";
    }
    return 0;
}
// Output: Acquired -> Released -> Caught
\end{lstlisting}

\begin{center}\rule{0.5\linewidth}{0.5pt}\end{center}

\hypertarget{debugging-strategies}{%
\section{3. DEBUGGING STRATEGIES}\label{debugging-strategies}}

\hypertarget{debug-macros-c98-style}{%
\subsection{3.1 Debug Macros (C++98
Style)}\label{debug-macros-c98-style}}

\begin{lstlisting}[language={C++}]
#include <iostream>

#ifdef DEBUG
    #define LOG(x) std::cout << "DEBUG: " << x << "\n"
#else
    #define LOG(x) 
#endif

int main() {
    int x = 10;
    LOG("x is " << x);
    return 0;
}
\end{lstlisting}

\hypertarget{professional-notes-chapter-72-exceptions}{%
\chapter{Professional Notes: Chapter 72:
Exceptions}\label{professional-notes-chapter-72-exceptions}}

Section 72.1: Catching exceptions Section 72.2: Rethrow (propagate)
exception Section 72.3: Best practice: throw by value, catch by const
reference Section 72.4: Custom exception Section 72.5:
std::uncaught\_exceptions Section 72.6: Function Try Block for regular
function Section 72.7: Nested exception Section 72.8: Function Try
Blocks In constructor Section 72.9: Function Try Blocks In destructor

\hypertarget{professional-notes-chapter-72-exceptions-1}{%
\chapter{Professional Notes: Chapter 72:
Exceptions}\label{professional-notes-chapter-72-exceptions-1}}

Section 72.1: Catching exceptions A try/catch block is used to catch
exceptions. The code in the try section is the code that may throw an
exception, and the code in the catch clause(s) handles the exception.
\#include \#include \#include int main() \{ std::string str(``foo'');
try \{ str.at(10); // access element, may throw std::out\_of\_range \}
catch (const std::out\_of\_range\& e) \{ // what() is inherited from
std::exception and contains an explanatory message std::cout
\textless\textless{} e.what(); \} \} Multiple catch clauses may be used
to handle multiple exception types. If multiple catch clauses are
present, the exception handling mechanism tries to match them in order
of their appearance in the code: std::string str(``foo''); try \{
str.reserve(2); // reserve extra capacity, may throw std::length\_error
str.at(10); // access element, may throw std::out\_of\_range \} catch
(const std::length\_error\& e) \{ std::cout \textless\textless{}
e.what(); \} catch (const std::out\_of\_range\& e) \{ std::cout
\textless\textless{} e.what(); \} Exception classes which are derived
from a common base class can be caught with a single catch clause for
the common base class. The above example can replace the two catch
clauses for std::length\_error and std::out\_of\_range with a single
clause for std:exception: std::string str(``foo''); try \{
str.reserve(2); // reserve extra capacity, may throw std::length\_error
str.at(10); // access element, may throw std::out\_of\_range \} catch
(const std::exception\& e) \{ std::cout \textless\textless{} e.what();
\} Because the catch clauses are tried in order, be sure to write more
specic catch clauses rst, otherwise your exception handling code might
never get called: try \{ /* Code throwing exceptions omitted. \emph{/ \}
catch (const std::exception\& e) \{ /} Handle all exceptions of type
std::exception. \emph{/ \} catch (const std::runtime\_error\& e) \{  /}
This block of code will never execute, because std::runtime\_error
inherits from std::exception, and all exceptions of type std::exception
were already caught by the previous catch clause. \emph{/ \} Another
possibility is the catch-all handler, which will catch any thrown
object: try \{ throw 10; \} catch (\ldots) \{ std::cout
\textless\textless{} ``caught an exception''; \} Section 72.2: Rethrow
(propagate) exception Sometimes you want to do something with the
exception you catch (like write to log or print a warning) and let it
bubble up to the upper scope to be handled. To do so, you can rethrow
any exception you catch: try \{ \ldots{} // some code here \} catch
(const SomeException\& e) \{ std::cout \textless\textless{} ``caught an
exception''; throw; \} Using throw; without arguments will re-throw the
currently caught exception. Version C++11 To rethrow a managed
std::exception\_ptr, the C++ Standard Library has the rethrow\_exception
function that can be used by including the header in your program.
\#include \#include \#include \#include void
handle\_eptr(std::exception\_ptr eptr) // passing by value is ok \{ try
\{ if (eptr) \{ std::rethrow\_exception(eptr); \} \} catch(const
std::exception\& e) \{ std::cout \textless\textless{} ``Caught exception
"'' \textless\textless{} e.what() \textless\textless{} ``"\n''; \} \}
int main() \{ std::exception\_ptr eptr; try \{ std::string().at(1); //
this generates an std::out\_of\_range \} catch(\ldots) \{ eptr =
std::current\_exception(); // capture \} handle\_eptr(eptr); \} //
destructor for std::out\_of\_range called here, when the eptr is
destructed Section 72.3: Best practice: throw by value, catch by const
reference In general, it is considered good practice to throw by value
(rather than by pointer), but catch by (const) reference. try \{ //
throw new std::runtime\_error(``Error!''); // Don't do this! // This
creates an exception object // on the heap and would require you to
catch the // pointer and manage the memory yourself. This can // cause
memory leaks! throw std::runtime\_error(``Error!''); \} catch (const
std::runtime\_error\& e) \{ std::cout \textless\textless{} e.what()
\textless\textless{} std::endl; \} One reason why catching by reference
is a good practice is that it eliminates the need to reconstruct the
object when being passed to the catch block (or when propagating through
to other catch blocks). Catching by reference also allows the exceptions
to be handled polymorphically and avoids object slicing. However, if you
are rethrowing an exception (like throw e;, see example below), you can
still get object slicing because the throw e; statement makes a copy of
the exception as whatever type is declared: \#include struct
BaseException \{ virtual const char} what() const \{ return
``BaseException''; \} \}; struct DerivedException : BaseException \{ //
``virtual'' keyword is optional here virtual const char* what() const \{
return ``DerivedException''; \} \}; int main(int argc, char** argv) \{
try \{ try \{ throw DerivedException(); \} catch (const BaseException\&
e) \{ std::cout \textless\textless{} ``First catch block:''
\textless\textless{} e.what() \textless\textless{} std::endl; // Output
==\textgreater{} First catch block: DerivedException throw e; // This
changes the exception to BaseException // instead of the original
DerivedException! \} \} catch (const BaseException\& e) \{ std::cout
\textless\textless{} ``Second catch block:'' \textless\textless{}
e.what() \textless\textless{} std::endl; // Output ==\textgreater{}
Second catch block: BaseException \} return 0; \} If you are sure that
you are not going to do anything to change the exception (like add
information or modify the message), catching by const reference allows
the compiler to make optimizations and can improve performance. But this
can still cause object splicing (as seen in the example above). Warning:
Beware of throwing unintended exceptions in catch blocks, especially
related to allocating extra memory or resources. For example,
constructing logic\_error, runtime\_error or their subclasses might
throw bad\_alloc due to memory running out when copying the exception
string, I/O streams might throw during logging with respective exception
masks set, etc. Section 72.4: Custom exception You shouldn't throw raw
values as exceptions, instead use one of the standard exception classes
or make your own. Having your own exception class inherited from
std::exception is a good way to go about it. Here's a custom exception
class which directly inherits from std::exception: \#include class
Except: virtual public std::exception \{ protected: int error\_number;
///\textless{} Error number int error\_offset; ///\textless{} Error
offset std::string error\_message; ///\textless{} Error message public:
/** Constructor (C++ STL string, int, int). * @param msg The error
message * @param err\_num Error number * @param err\_off Error offset */
explicit Except(const std::string\& msg, int err\_num, int err\_off):
error\_number(err\_num), error\_offset(err\_off), error\_message(msg)
\{\} /** Destructor. * Virtual to allow for subclassing. \emph{/ virtual
\textasciitilde Except() throw () \{\} /\textbf{ Returns a pointer to
the (constant) error description. * @return A pointer to a const
char\emph{. The underlying memory } is in possession of the Except
object. Callers must * not attempt to free the memory. \emph{/ virtual
const char} what() const throw () \{ return error\_message.c\_str(); \}
/} Returns error number. } @return \#error\_number */ virtual int
getErrorNumber() const throw() \{ return error\_number; \} /**Returns
error offset. * @return \#error\_offset */ virtual int getErrorOffset()
const throw() \{  return error\_offset; \} \}; An example throw catch:
try \{ throw(Except(``Couldn't do what you were expecting'', -12, -34));
\} catch (const Except\& e) \{ std::cout\textless\textless e.what()
\textless\textless{}``\nError number:''\textless\textless e.getErrorNumber()
\textless\textless{}``\nError offset:''\textless\textless e.getErrorOffset();
\} As you are not only just throwing a dumb error message, also some
other values representing what the error exactly was, your error
handling becomes much more ecient and meaningful. There's an exception
class that let's you handle error messages nicely :std::runtime\_error
You can inherit from this class too: \#include class Except: virtual
public std::runtime\_error \{ protected: int error\_number;
///\textless{} Error number int error\_offset; ///\textless{} Error
offset public: /** Constructor (C++ STL string, int, int). * @param msg
The error message * @param err\_num Error number * @param err\_off Error
offset */ explicit Except(const std::string\& msg, int err\_num, int
err\_off): std::runtime\_error(msg) \{ error\_number = err\_num;
error\_offset = err\_off; \} /** Destructor. * Virtual to allow for
subclassing. */ virtual \textasciitilde Except() throw () \{\} /**
Returns error number. * @return \#error\_number */ virtual int
getErrorNumber() const throw() \{ return error\_number; \} /**Returns
error offset.  * @return \#error\_offset */ virtual int
getErrorOffset() const throw() \{ return error\_offset; \} \}; Note that
I haven't overridden the what() function from the base class
(std::runtime\_error) i.e we will be using the base class's version of
what(). You can override it if you have further agenda. Section 72.5:
std::uncaught\_exceptions Version c++17 C++17 introduces int
std::uncaught\_exceptions() (to replace the limited bool
std::uncaught\_exception()) to know how many exceptions are currently
uncaught. That allows for a class to determine if it is destroyed during
a stack unwinding or not. \#include \#include \#include // Apply change
on destruction: // Rollback in case of exception (failure) // Else
Commit (success) class Transaction \{ public: Transaction(const
std::string\& s) : message(s) \{\} Transaction(const Transaction\&) =
delete; Transaction\& operator =(const Transaction\&) = delete; void
Commit() \{ std::cout \textless\textless{} message \textless\textless{}
``: Commit\n''; \} void RollBack() noexcept(true) \{ std::cout
\textless\textless{} message \textless\textless{} ``: Rollback\n''; \}
// \ldots{} \textasciitilde Transaction() \{ if (uncaughtExceptionCount
== std::uncaught\_exceptions()) \{ Commit(); // May throw. \} else \{ //
current stack unwinding RollBack(); \} \} private: std::string message;
int uncaughtExceptionCount = std::uncaught\_exceptions(); \}; class Foo
\{ public: \textasciitilde Foo() \{ try \{ Transaction transaction(``In
\textasciitilde Foo''); // Commit, // even if there is an uncaught
exception //\ldots{} \} catch (const std::exception\& e) \{ std::cerr
\textless\textless{} ``exception/\textasciitilde Foo:''
\textless\textless{} e.what() \textless\textless{} std::endl;  \} \}
\}; int main() \{ try \{ Transaction transaction(``In main''); //
RollBack Foo foo; // \textasciitilde Foo commit its transaction.
//\ldots{} throw std::runtime\_error(``Error''); \} catch (const
std::exception\& e) \{ std::cerr \textless\textless{}
``exception/main:'' \textless\textless{} e.what() \textless\textless{}
std::endl; \} \} Output: In \textasciitilde Foo: Commit In main:
Rollback exception/main:Error Section 72.6: Function Try Block for
regular function void function\_with\_try\_block() try \{ // try block
body \} catch (\ldots) \{ // catch block body \} Which is equivalent to
void function\_with\_try\_block() \{ try \{ // try block body \} catch
(\ldots) \{ // catch block body \} \} Note that for constructors and
destructors, the behavior is dierent as the catch block re-throws an
exception anyway (the caught one if there is no other throw in the catch
block body). The function main is allowed to have a function try block
like any other function, but main's function try block will not catch
exceptions that occur during the construction of a non-local static
variable or the destruction of any static variable. Instead,
std::terminate is called. Section 72.7: Nested exception Version C++11
During exception handling there is a common use case when you catch a
generic exception from a low-level function (such as a lesystem error or
data transfer error) and throw a more specic high-level exception which
indicates that some high-level operation could not be performed (such as
being unable to publish a photo on Web). This allows exception handling
to react to specic problems with high level operations and also allows,
having only error an message, the programmer to nd a place in the
application where an exception occurred. Downside of this solution is
that exception callstack is truncated and original exception is lost.
This forces developers to manually include text of original exception
into a newly created one. Nested exceptions aim to solve the problem by
attaching low-level exception, which describes the cause, to a high
level exception, which describes what it means in this particular case.
std::nested\_exception allows to nest exceptions thanks to
std::throw\_with\_nested: \#include \#include \#include \#include
\#include struct MyException \{ MyException(const std::string\& message)
: message(message) \{\} std::string message; \}; void
print\_current\_exception(int level) \{ try \{ throw; \} catch (const
std::exception\& e) \{ std::cerr \textless\textless{} std::string(level,
' `) \textless\textless{} ``exception:'' \textless\textless{} e.what()
\textless\textless{}'\n`; \} catch (const MyException\& e) \{ std::cerr
\textless\textless{} std::string(level,' `) \textless\textless{}
``MyException:'' \textless\textless{} e.message
\textless\textless{}'\n'; \} catch (\ldots) \{ std::cerr
\textless\textless{} ``Unkown exception\n''; \} \} void
print\_current\_exception\_with\_nested(int level = 0) \{ try \{ throw;
\} catch (\ldots) \{ print\_current\_exception(level); \}\\
try \{ throw; \} catch (const std::nested\_exception\& nested) \{ try \{
nested.rethrow\_nested(); \} catch (\ldots) \{
print\_current\_exception\_with\_nested(level + 1); // recursion \} \}
catch (\ldots) \{ //Empty // End recursion \} \} // sample function that
catches an exception and wraps it in a nested exception void
open\_file(const std::string\& s) \{ try \{ std::ifstream file(s);
file.exceptions(std::ios\_base::failbit); \} catch(\ldots) \{
std::throw\_with\_nested(MyException\{``Couldn't open'' + s\}); \} \} //
sample function that catches an exception and wraps it in a nested
exception void run() \{ try \{ open\_file(``nonexistent.file''); \}
catch(\ldots) \{ std::throw\_with\_nested( std::runtime\_error(``run()
failed'') ); \} \} // runs the sample function above and prints the
caught exception int main() \{ try \{ run(); \} catch(\ldots) \{
print\_current\_exception\_with\_nested(); \} \} Possible output:
exception: run() failed MyException: Couldn't open nonexistent.file
exception: basic\_ios::clear If you work only with exceptions inherited
from std::exception, code can even be simplied. Section 72.8: Function
Try Blocks In constructor The only way to catch exception in initializer
list: struct A : public B \{ A() try : B(), foo(1), bar(2) \{ //
constructor body \} catch (\ldots) \{ // exceptions from the initializer
list and constructor are caught here // if no exception is thrown here
// then the caught exception is re-thrown. \} private: Foo foo; Bar bar;
\}; Section 72.9: Function Try Blocks In destructor struct A \{
\textasciitilde A() noexcept(false) try \{ // destructor body \} catch
(\ldots) \{ // exceptions of destructor body are caught here // if no
exception is thrown here // then the caught exception is re-thrown. \}
\}; Note that, although this is possible, one needs to be very careful
with throwing from destructor, as if a destructor called during stack
unwinding throws an exception, std::terminate is called.

\begin{center}\rule{0.5\linewidth}{0.5pt}\end{center}

\hypertarget{chapter-10-advanced-streams-file-io}{%
\chapter{Chapter 10: Advanced Streams \& File
I/O}\label{chapter-10-advanced-streams-file-io}}

This chapter explores the full power of the C++
\passthrough{\lstinline!<iostream>!},
\passthrough{\lstinline!<fstream>!}, and
\passthrough{\lstinline!<sstream>!} libraries, deconstructing how they
interact with the filesystem and formatted data.

\hypertarget{file-io-foundations}{%
\section{10.1 File I/O Foundations}\label{file-io-foundations}}

Working with files involves the \passthrough{\lstinline!std::fstream!},
\passthrough{\lstinline!std::ifstream!}, and
\passthrough{\lstinline!std::ofstream!} classes.

\hypertarget{opening-modes}{%
\subsection{1. Opening Modes}\label{opening-modes}}

\begin{itemize}
\tightlist
\item
  \passthrough{\lstinline!std::ios::in!}: Open for reading.
\item
  \passthrough{\lstinline!std::ios::out!}: Open for writing (overwrites
  existing).
\item
  \passthrough{\lstinline!std::ios::app!}: Append to end of file.
\item
  \passthrough{\lstinline!std::ios::ate!}: Open and seek to end.
\item
  \passthrough{\lstinline!std::ios::binary!}: Binary mode (no CRLF
  translation).
\end{itemize}

\hypertarget{binary-vs-text-mode}{%
\subsection{2. Binary vs Text Mode}\label{binary-vs-text-mode}}

In text mode, special characters like \passthrough{\lstinline!\\n!}
might be translated (e.g., to \passthrough{\lstinline!\\r\\n!} on
Windows). Binary mode ensures that exactly what you write is what
appears on disk.

\begin{lstlisting}[language={C++}]
std::ofstream file("data.bin", std::ios::binary);
double d = 3.14;
file.write(reinterpret_cast<const char*>(&d), sizeof(d));
\end{lstlisting}

\begin{center}\rule{0.5\linewidth}{0.5pt}\end{center}

\hypertarget{stream-manipulators}{%
\section{10.2 Stream Manipulators}\label{stream-manipulators}}

Manipulators allow you to change how data is formatted on the fly.

\hypertarget{numeric-formatting}{%
\subsection{1. Numeric Formatting}\label{numeric-formatting}}

\begin{itemize}
\tightlist
\item
  \passthrough{\lstinline!std::hex!},
  \passthrough{\lstinline!std::oct!},
  \passthrough{\lstinline!std::dec!}: Change base.
\item
  \passthrough{\lstinline!std::setprecision(n)!}: Set floating point
  precision (requires \passthrough{\lstinline!<iomanip>!}).
\item
  \passthrough{\lstinline!std::fixed!},
  \passthrough{\lstinline!std::scientific!}: Change notation.
\end{itemize}

\hypertarget{padding-and-alignment}{%
\subsection{2. Padding and Alignment}\label{padding-and-alignment}}

\begin{itemize}
\tightlist
\item
  \passthrough{\lstinline!std::setw(n)!}: Set width of next field.
\item
  \passthrough{\lstinline!std::setfill(c)!}: Set fill character.
\item
  \passthrough{\lstinline!std::left!},
  \passthrough{\lstinline!std::right!},
  \passthrough{\lstinline!std::internal!}: Change alignment.
\end{itemize}

\begin{center}\rule{0.5\linewidth}{0.5pt}\end{center}

\hypertarget{string-streams-sstream}{%
\section{\texorpdfstring{10.3 String Streams
(\texttt{sstream})}{10.3 String Streams (sstream)}}\label{string-streams-sstream}}

\passthrough{\lstinline!std::stringstream!} allows you to treat a string
like a stream, enabling easy conversion between types and strings.

\begin{lstlisting}[language={C++}]
#include <sstream>
std::stringstream ss;
ss << "The answer is " << 42;
std::string s = ss.str();
\end{lstlisting}

\begin{center}\rule{0.5\linewidth}{0.5pt}\end{center}

\hypertarget{professional-notes-stream-architecture}{%
\subsection{Professional Notes: Stream
Architecture}\label{professional-notes-stream-architecture}}

\hypertarget{the-ios_base-state-machine}{%
\subsubsection{\texorpdfstring{1. The \texttt{ios\_base} State
Machine}{1. The ios\_base State Machine}}\label{the-ios_base-state-machine}}

Every stream inherits from \passthrough{\lstinline!ios\_base!}, which
maintains internal flags for formatting and errors. *
\textbf{Performance Trap}: Streams are significantly slower than C's
\passthrough{\lstinline!printf!}/\passthrough{\lstinline!scanf!} due to
the overhead of object construction and virtual function calls. *
\textbf{Optimization}:
\passthrough{\lstinline!std::ios::sync\_with\_stdio(false);!} disables
the synchronization between C++ and C streams, making
\passthrough{\lstinline!std::cin!} as fast as
\passthrough{\lstinline!scanf!}.

\hypertarget{stream-buffering-streambuf}{%
\subsubsection{\texorpdfstring{2. Stream Buffering
(\texttt{streambuf})}{2. Stream Buffering (streambuf)}}\label{stream-buffering-streambuf}}

The actual data transfer is handled by a buffer object
(\passthrough{\lstinline!rdbuf!}). * \textbf{Redirection}: You can
redirect \passthrough{\lstinline!cout!} to a file by swapping its
buffer:

\begin{lstlisting}[language={C++}]
std::ofstream out("log.txt");
std::streambuf *coutbuf = std::cout.rdbuf(); 
std::cout.rdbuf(out.rdbuf()); // cout now writes to log.txt
\end{lstlisting}

\hypertarget{custom-manipulators}{%
\subsubsection{3. Custom Manipulators}\label{custom-manipulators}}

You can create your own manipulators by defining functions that take and
return a reference to a stream:

\begin{lstlisting}[language={C++}]
std::ostream& tab(std::ostream& os) {
    return os << "\t";
}
std::cout << "Data1" << tab << "Data2" << std::endl;
\end{lstlisting}

\begin{center}\rule{0.5\linewidth}{0.5pt}\end{center}

\begin{center}\rule{0.5\linewidth}{0.5pt}\end{center}

\hypertarget{volume-02-modern-revolution-c11}{%
\chapter{VOLUME 02 MODERN REVOLUTION
C11}\label{volume-02-modern-revolution-c11}}
